\chapter*{A Note on the Title}
\addcontentsline{toc}{chapter}{A Note on the Title}
\markboth{A Note on the Title}{A Note on the Title}

The word \emph{hypernormal} is borrowed, with some distortion, from Alexei Yurchak's account of the last decades of the Soviet Union~\cite{yurchak2006}, in which the forms of public life became so fixed and so faithfully reproduced that nobody could say what they were for, and from Adam Curtis's use of the term as the title of a documentary~\cite{curtis2016}. This book is not about either. It is about counting, and it uses the word for one narrow property of an edition. The book does not argue that hypernormality is a cultural condition or a description of contemporary media, and it does not contest the literature the word comes from. The two concerns meet at one point only: when editions differ, what viewers have in common to talk about is thinner than it looks, as the occasion for talk of Section~1.1.2 and the worldwide premiere of Chapter~11 show, and the book counts how much thinner.

An edition is \emph{hypernormal} here when it satisfies every rule in force and shows the viewer no seam of its assembly: nothing in it reads as a variant, a splice, or a substitution. What the story shows is a separate matter. A character may miss a strip of lawn, and that mistake is part of the film, something it may need the viewer to see. The grass that was mown and then stands uncut again is the other kind, a defect of assembly that a viewer may notice and report. The lawn scene of Chapter~1 is the standard case. Each viewer sees a lawn mown in a plausible order, with the machine travelling in a consistent direction across the cuts, and nothing about it looks like a variant. Across the audience, however, no two viewers hold the same edition, and the production has chosen on purpose to keep it so. Normality is therefore a property that each edition has on its own, and sameness is a property that no viewer can check, because checking requires access to someone else's film.

The title joins this to the second half of the question the book asks, which is how much variation a production can sustain while remaining open to correction. Correction in these chapters begins with a report, and a report begins with a viewer who notices a defect of assembly. A hypernormal edition hides its seams from every viewer, so the errors that matter most are the ones the rules did not anticipate and that each viewer sees as ordinary. A viewer who notices a character's mistake has noticed the film, and a viewer who notices the grass has noticed the assembly, and only the second kind of noticing starts a correction. The ledger of Chapter~12 and the adjudication of Chapter~7 let a report, once it is made, be traced to the viewers who hold an affected edition. They do not make an unnoticed defect noticed. That limit is what the second half of the title names, and Chapter~18 treats how corrections arrive as an empirical question the mathematics leaves open.

Two qualifications apply. The word describes the models of the chapters and says nothing about any real audience, since no experiment in this book measures what viewers notice. And the satire, to the extent there is any, is meant to be carried by the arithmetic: the counts are exact on the instances given, and the reader is invited to decide what they show.
